\documentclass[11pt]{article}
\usepackage[utf8]{inputenc}
\usepackage{amsmath, amssymb, amsthm}
\usepackage{algorithm}
\usepackage{algorithmic}
\usepackage{geometry}
\geometry{margin=1in}

\title{Analysis of Online Gradient Descent (OGD)}
\author{}
\date{}

\begin{document}

\maketitle

\section*{Algorithm Name}
Online Gradient Descent (OGD)

\section*{Mathematical Formulation}
Online Gradient Descent is designed for online learning scenarios where an algorithm must make sequential decisions and subsequently receives feedback in the form of a loss function. The goal is to minimize the \textit{regret}, which is the difference between the cumulative loss of the algorithm and the cumulative loss of the best fixed decision in hindsight.

Let $\mathcal{W} \subset \mathbb{R}^d$ be a convex, closed, and bounded feasible set. At each time step $t = 1, 2, \dots, T$:
\begin{enumerate}
    \item The algorithm predicts a parameter vector $w_t \in \mathcal{W}$.
    \item The environment reveals a convex loss function $f_t : \mathcal{W} \to \mathbb{R}$.
    \item The algorithm suffers loss $f_t(w_t)$.
    \item The algorithm computes the gradient $g_t = \nabla f_t(w_t)$ (or a subgradient if $f_t$ is not differentiable).
    \item The parameters are updated using the projected gradient descent rule:
    $$w_{t+1} = \Pi_{\mathcal{W}}(w_t - \eta_t g_t)$$
    where $\eta_t > 0$ is the learning rate (step size) and $\Pi_{\mathcal{W}}$ is the Euclidean projection onto the set $\mathcal{W}$, defined as $\Pi_{\mathcal{W}}(v) = \arg\min_{w \in \mathcal{W}} \|w - v\|_2$.
\end{enumerate}

\section*{Pseudocode}
\begin{algorithm}[H]
\caption{Online Gradient Descent}
\begin{algorithmic}[1]
\REQUIRE Feasible set $\mathcal{W}$, number of iterations $T$, learning rate schedule $\{\eta_t\}_{t=1}^T$
\STATE Initialize $w_1 \in \mathcal{W}$ arbitrarily (e.g., $w_1 = \mathbf{0}$)
\FOR{$t = 1$ \TO $T$}
    \STATE Output prediction $w_t$
    \STATE Receive online loss function $f_t$ and suffer loss $f_t(w_t)$
    \STATE Compute gradient $g_t = \nabla f_t(w_t)$
    \STATE Update parameters: $w_{t+1} = \Pi_{\mathcal{W}}(w_t - \eta_t g_t)$
\ENDFOR
\end{algorithmic}
\end{algorithm}

\section*{Dry Run Example}
Let the feasible set be $\mathcal{W} = \mathbb{R}$ (so projection is the identity). Let the online loss function be identical at each step for simplicity: $f_t(w) = (w-3)^2$. We initialize $w_1 = 0$ and use a constant learning rate $\eta = 0.1$. We run for $T=3$ iterations.

\textbf{Iteration 1 ($t=1$):}
\begin{itemize}
    \item Predict: $w_1 = 0$
    \item Receive loss function: $f_1(w) = (w-3)^2$
    \item Suffer loss: $f_1(w_1) = (0-3)^2 = 9$
    \item Compute gradient: $g_1 = \nabla f_1(w_1) = 2(w_1 - 3) = 2(0 - 3) = -6$
    \item Update: $w_2 = w_1 - \eta g_1 = 0 - 0.1(-6) = 0.6$
\end{itemize}

\textbf{Iteration 2 ($t=2$):}
\begin{itemize}
    \item Predict: $w_2 = 0.6$
    \item Receive loss function: $f_2(w) = (w-3)^2$
    \item Suffer loss: $f_2(w_2) = (0.6-3)^2 = (-2.4)^2 = 5.76$
    \item Compute gradient: $g_2 = \nabla f_2(w_2) = 2(0.6 - 3) = -4.8$
    \item Update: $w_3 = w_2 - \eta g_2 = 0.6 - 0.1(-4.8) = 1.08$
\end{itemize}

\textbf{Iteration 3 ($t=3$):}
\begin{itemize}
    \item Predict: $w_3 = 1.08$
    \item Receive loss function: $f_3(w) = (w-3)^2$
    \item Suffer loss: $f_3(w_3) = (1.08-3)^2 = (-1.92)^2 \approx 3.6864$
    \item Compute gradient: $g_3 = \nabla f_3(w_3) = 2(1.08 - 3) = -3.84$
    \item Update: $w_4 = w_3 - \eta g_3 = 1.08 - 0.1(-3.84) = 1.464$
\end{itemize}
Notice how the loss monotonically decreases as the parameter approaches the optimum at $w=3$.

\section*{Assumptions}
To prove the sublinear regret bound, we require the following formal assumptions:
\begin{enumerate}
    \item \textbf{Convexity of Domain:} The decision set $\mathcal{W} \subset \mathbb{R}^d$ is closed, convex, and non-empty.
    \item \textbf{Bounded Domain:} The diameter of $\mathcal{W}$ is bounded by $D$. That is, for all $x, y \in \mathcal{W}$, $\|x - y\|_2 \le D$.
    \item \textbf{Convexity of Losses:} For all $t$, the loss function $f_t: \mathcal{W} \to \mathbb{R}$ is convex.
    \item \textbf{Lipschitz Continuous Losses (Bounded Gradients):} For all $t$, the function $f_t$ is $G$-Lipschitz over $\mathcal{W}$. This implies that for all $w \in \mathcal{W}$, the gradient is bounded: $\|\nabla f_t(w)\|_2 \le G$.
\end{enumerate}

\section*{Main Convergence Theorem}
\textbf{Theorem (Regret Bound of OGD):} Under Assumptions 1-4, if Online Gradient Descent is run with a constant learning rate $\eta_t = \eta = \frac{D}{G\sqrt{T}}$, then the regret $R_T$ after $T$ steps is bounded by:
$$R_T = \sum_{t=1}^T f_t(w_t) - \min_{w^* \in \mathcal{W}} \sum_{t=1}^T f_t(w^*) \le D G \sqrt{T}$$
Since $R_T \in \mathcal{O}(\sqrt{T})$, the average regret $\frac{R_T}{T} \in \mathcal{O}(1/\sqrt{T})$, which approaches $0$ as $T \to \infty$.

\section*{Theoretical Properties}
\begin{itemize}
    \item \textbf{Sublinear Regret:} The $\mathcal{O}(\sqrt{T})$ regret bound shows that in the long run, the algorithm performs almost as well as the best fixed decision in hindsight.
    \item \textbf{Hyperparameter Sensitivity:} The optimal fixed learning rate requires knowing the time horizon $T$ in advance. If $T$ is unknown, a time-varying learning rate $\eta_t = \frac{D}{G\sqrt{t}}$ can be used, which yields a slightly worse but still sublinear regret bound of $\mathcal{O}(\sqrt{T})$.
    \item \textbf{Comparison to Baselines:} Unlike standard Gradient Descent (GD) which assumes a fixed objective function $f(w)$, OGD handles a sequence of potentially adversarial functions $f_t(w)$, providing a much stronger robustness guarantee.
\end{itemize}

\section*{Discussion}
The theoretical framework of OGD bridges optimization and online learning. The upper bound on regret is tight for convex and Lipschitz functions; one cannot achieve a generally better rate than $\Omega(\sqrt{T})$ without further assumptions (like strong convexity).

\section*{Preliminaries (Definitions and Lemmas)}
\textbf{Definition 1 (Regret):} The regret relative to a fixed optimal parameter $w^* = \arg\min_{w \in \mathcal{W}} \sum_{t=1}^T f_t(w)$ is defined as:
$$R_T = \sum_{t=1}^T f_t(w_t) - \sum_{t=1}^T f_t(w^*)$$

\textbf{Lemma 1 (Properties of Projection):} For any convex, closed set $\mathcal{W}$, and any $v \in \mathbb{R}^d, w \in \mathcal{W}$, the Euclidean projection satisfies:
$$\|\Pi_{\mathcal{W}}(v) - w\|^2 \le \|v - w\|^2$$

\textbf{Lemma 2 (First-Order Convexity Condition):} For a differentiable convex function $f$, for any $x, y$:
$$f(x) - f(y) \le \langle \nabla f(x), x - y \rangle$$

\section*{Main Proof (Step-by-Step Derivation)}
Let $w^* \in \mathcal{W}$ be the optimal fixed decision in hindsight. We track the distance between our current parameter $w_t$ and $w^*$.

\begin{align*}
\text{[Step 1]} \quad & \|w_{t+1} - w^*\|^2 = \|\Pi_{\mathcal{W}}(w_t - \eta_t g_t) - w^*\|^2 \\
& \text{Substitute the update rule defining } w_{t+1}. \\
\text{[Step 2]} \quad & \|w_{t+1} - w^*\|^2 \le \|w_t - \eta_t g_t - w^*\|^2 \\
& \text{Apply Lemma 1 (Projection Property) since } w^* \in \mathcal{W}. \\
\text{[Step 3]} \quad & \|w_{t+1} - w^*\|^2 \le \|w_t - w^*\|^2 - 2\eta_t \langle g_t, w_t - w^* \rangle + \eta_t^2 \|g_t\|^2 \\
& \text{Expand the squared Euclidean norm } \|a - b\|^2 = \|a\|^2 - 2\langle a, b \rangle + \|b\|^2. \\
\text{[Step 4]} \quad & 2\eta_t \langle g_t, w_t - w^* \rangle \le \|w_t - w^*\|^2 - \|w_{t+1} - w^*\|^2 + \eta_t^2 \|g_t\|^2 \\
& \text{Rearrange the inequality to isolate the inner product term.} \\
\text{[Step 5]} \quad & \langle g_t, w_t - w^* \rangle \le \frac{\|w_t - w^*\|^2 - \|w_{t+1} - w^*\|^2}{2\eta_t} + \frac{\eta_t}{2} \|g_t\|^2 \\
& \text{Divide the entire inequality by } 2\eta_t. \\
\text{[Step 6]} \quad & \langle g_t, w_t - w^* \rangle \le \frac{\|w_t - w^*\|^2 - \|w_{t+1} - w^*\|^2}{2\eta_t} + \frac{\eta_t}{2} G^2 \\
& \text{Apply Assumption 4 (Lipschitz gradients), substituting } \|g_t\|^2 \le G^2. \\
\text{[Step 7]} \quad & f_t(w_t) - f_t(w^*) \le \langle g_t, w_t - w^* \rangle \\
& \text{Apply Lemma 2 (Convexity) substituting } x=w_t, y=w^*, \text{ and } \nabla f_t(w_t)=g_t. \\
\text{[Step 8]} \quad & f_t(w_t) - f_t(w^*) \le \frac{\|w_t - w^*\|^2 - \|w_{t+1} - w^*\|^2}{2\eta_t} + \frac{\eta_t}{2} G^2 \\
& \text{Chain the inequalities from Step 6 and Step 7.} \\
\text{[Step 9]} \quad & \sum_{t=1}^T \left( f_t(w_t) - f_t(w^*) \right) \le \sum_{t=1}^T \left( \frac{\|w_t - w^*\|^2 - \|w_{t+1} - w^*\|^2}{2\eta_t} + \frac{\eta_t}{2} G^2 \right) \\
& \text{Sum the inequality over all time steps } t=1 \text{ to } T \text{ to bound the total regret.}
\end{align*}

\section*{Rate Derivation (With Constants)}
We now simplify the bounded sum, assuming a constant learning rate $\eta_t = \eta$ for all $t$.

\begin{align*}
\text{[Step 10]} \quad & R_T \le \frac{1}{2\eta} \sum_{t=1}^T \left( \|w_t - w^*\|^2 - \|w_{t+1} - w^*\|^2 \right) + \sum_{t=1}^T \frac{\eta G^2}{2} \\
& \text{Factor out } \frac{1}{2\eta} \text{ and separate the sums.} \\
\text{[Step 11]} \quad & R_T \le \frac{1}{2\eta} \left( \|w_1 - w^*\|^2 - \|w_{T+1} - w^*\|^2 \right) + \frac{\eta G^2 T}{2} \\
& \text{Collapse the telescoping sum.} \\
\text{[Step 12]} \quad & R_T \le \frac{1}{2\eta} \|w_1 - w^*\|^2 + \frac{\eta G^2 T}{2} \\
& \text{Drop the non-positive term } -\|w_{T+1} - w^*\|^2 \le 0. \\
\text{[Step 13]} \quad & R_T \le \frac{D^2}{2\eta} + \frac{\eta G^2 T}{2} \\
& \text{Apply Assumption 2 (Bounded Domain), substituting } \|w_1 - w^*\|^2 \le D^2. \\
\text{[Step 14]} \quad & \frac{d}{d\eta} \left( \frac{D^2}{2\eta} + \frac{\eta G^2 T}{2} \right) = -\frac{D^2}{2\eta^2} + \frac{G^2 T}{2} = 0 \\
& \text{Differentiate the upper bound with respect to } \eta \text{ and set to zero to find the optimal } \eta. \\
\text{[Step 15]} \quad & \eta^* = \frac{D}{G\sqrt{T}} \\
& \text{Solve for the optimal constant learning rate } \eta^*. \\
\text{[Step 16]} \quad & R_T \le \frac{D^2}{2 \left( \frac{D}{G\sqrt{T}} \right)} + \frac{\left( \frac{D}{G\sqrt{T}} \right) G^2 T}{2} \\
& \text{Substitute } \eta^* \text{ back into the regret bound from Step 13.} \\
\text{[Step 17]} \quad & R_T \le \frac{D G \sqrt{T}}{2} + \frac{D G \sqrt{T}}{2} = D G \sqrt{T} \\
& \text{Simplify the algebraic expression to obtain the final sublinear regret bound.}
\end{align*}

\section*{Special Cases}
\textbf{Strongly Convex Functions:}
If the sequence of loss functions $f_t$ are uniformly $\mu$-strongly convex, the convexity inequality (Lemma 2) gets an additional term:
$$f_t(w_t) - f_t(w^*) \le \langle g_t, w_t - w^* \rangle - \frac{\mu}{2}\|w_t - w^*\|^2$$
By employing a time-decaying learning rate specifically tuned to the strong convexity parameter, $\eta_t = \frac{1}{\mu t}$, the regret bound drastically improves from $\mathcal{O}(\sqrt{T})$ to $\mathcal{O}(\log T)$. The proof follows a similar telescoping structure but leverages the negative quadratic term to cancel out the variance accumulation over time.

\end{document}